% ----------------------------------------------------------------------
%  3.4. Диффузия и кинетика фазовых превращений
%  3.4.1. Линейные феноменологические законы
% ----------------------------------------------------------------------
\begin{orientir}
Почти всё, что происходит со сплавом при нагреве,~— гомогенизация,
старение, науглероживание, рост интерметаллидных прослоек, спекание,
ползучесть,~— идёт через \emph{диффузию}: блуждание атомов по решётке.
Раздел построен от общего к частному: 3.4.1~— общие законы переноса и
их термодинамика (что на самом деле «движет» атомы); 3.4.2~— уравнения
Фика и их решения; 3.4.3~— как прыгают атомы и почему коэффициент
диффузии так сильно зависит от температуры и дефектов; 3.4.4~— эффект
Киркендалла и теория Даркена; 3.4.5~— диффузия при фазовых
превращениях. Многие величины здесь знакомы по главе~2 (вакансии,
энергия миграции).
\end{orientir}

\subsection{Линейные феноменологические законы}\label{ssec:onsager}

Явления можно описывать двумя дополняющими друг друга способами.
\emph{Молекулярно-кинетический} подход строит модель движения частиц.
\emph{Феноменологический}~— устанавливает общие закономерности, не
опираясь на модель строения вещества; в этом его сила. Важнейшая
феноменологическая теория~— термодинамика.

\paragraph{Потоки и силы}
Многие необратимые процессы подчиняются одинаковому закону: поток
пропорционален градиенту (табл.~\ref{tab:linear-laws}). Первым был
закон Фурье для теплопроводности; Фик перенёс его на диффузию:
\begin{keyeq}\label{eq:fick1-general}
\vb J_i=-D_i\operatorname{grad}C_i .
\end{keyeq}

\begin{table}[ht]
\centering
\caption{Линейные законы переноса}
\label{tab:linear-laws}
\small
\begin{tabular}{@{}llll@{}}
\toprule
Процесс & Поток & Причина («сила») & Закон \\
\midrule
теплопроводность & тепла $\vb q$ & $-\operatorname{grad}T$ & Фурье: $\vb q=-\varkappa\operatorname{grad}T$ \\
диффузия & вещества $\vb J$ & $-\operatorname{grad}C$ & Фик: $\vb J=-D\operatorname{grad}C$ \\
электропроводность & заряда $\vb j$ & $-\operatorname{grad}\varphi$ & Ом: $\vb j=-\sigma\operatorname{grad}\varphi$ \\
\bottomrule
\end{tabular}
\end{table}

Когда два процесса идут одновременно, они \emph{налагаются} и порождают
новые, перекрёстные эффекты:
\begin{itemize}
\item теплопроводность + электропроводность~— термоэлектрические
явления: эффекты Пельтье и Томсона (выделение или поглощение тепла на
спаях или по длине проводника при прохождении тока) и
термо-ЭДС в цепи со спаями при разной температуре;
\item диффузия + теплопроводность~— \emph{термодиффузия}: градиент
температуры создаёт градиент концентрации (эффект Соре), а градиент
концентрации~— градиент температуры (эффект Дюфура);
\item диффузия + электропроводность~— диффузионный потенциал;
\item в анизотропном кристалле градиент температуры в одном направлении
может вызвать поток тепла в другом.
\end{itemize}
Математически это описывают, добавляя к линейному закону ещё одно
слагаемое: например, при термодиффузии поток вещества пропорционален и
градиенту концентрации, и градиенту температуры.

\paragraph{Теория Онзагера}
Строгую термодинамику необратимых процессов построил Онзагер. Причины
необратимых явлений (градиенты температуры, концентрации, потенциала)
называют \term{силами} $X_k$, а вызванные ими процессы (потоки тепла,
вещества, заряда)~— \term{потоками} $J_i$. В общем случае любой поток
вызывается всеми силами:
\begin{keyeq}\label{eq:onsager}
J_i=\sum_{k=1}^nL_{ik}X_k ,\qquad L_{ik}=L_{ki}.
\end{keyeq}
Диагональные коэффициенты $L_{ii}$ описывают «обычные» явления
(диффузию, теплопроводность), недиагональные $L_{ik}$~— перекрёстные
(термодиффузию, эффект Дюфура). Второе равенство~— \term{соотношения
взаимности Онзагера}: при правильном выборе потоков и сил матрица
коэффициентов симметрична. Они отражают на макроуровне
\emph{микроскопическую обратимость}: уравнения движения частиц не
меняются при замене $t\to-t$, и в равновесии любой элементарный процесс
и обратный ему идут в среднем с одинаковой скоростью. Онзагер получил
их, рассматривая флуктуации в равновесной системе.

Как выбрать «правильные» потоки и силы? Их подсказывает
\term{производство энтропии}~— скорость её роста при необратимом
процессе:
\begin{equation}\label{eq:entropy-prod}
\sigma=\frac1T\frac{\partial(\Delta S)}{\partial t}=\frac1T\sum_iJ_iX_i\ \ge0 .
\end{equation}
Потоки~— производные по времени параметров состояния, силы~—
производные энтропии по этим параметрам.

\paragraph{Что на самом деле движет атомы}
Применим этот подход к изотермической диффузии.
\begin{shagi}
\item При постоянных $T$ и $p$ производство энтропии пропорционально
скорости убыли свободной энергии Гиббса: $\sigma=-\frac1T\frac{\d G}{\d t}$
(де Донде).
\item Перенос $n_i$ молей вещества из плоскости с химическим
потенциалом $\mu_i$ в плоскость с $\mu_i+\frac{\partial\mu_i}{\partial
x}\d x$ меняет $G$ на $n_i\frac{\partial\mu_i}{\partial x}\d x$.
\item Отсюда на единицу объёма
$\sigma=-\frac1T\sum_iJ_i\operatorname{grad}\mu_i$. Сравнивая с
\eqref{eq:entropy-prod}, находим силу: $X_i=-\operatorname{grad}\mu_i$.
\end{shagi}
\begin{keyeq}\label{eq:J-mu}
J_i=-\sum_{k=1}^nL_{ik}\operatorname{grad}\mu_k .
\end{keyeq}

\begin{glavnoe}
Движущая сила диффузии~— градиент \emph{химического потенциала}, а не
концентрации. Закон Фика~— частный случай: в идеальном растворе $\mu$
растёт с концентрацией монотонно, и оба градиента направлены
одинаково. Поэтому коэффициент $D$ в законе Фика, вообще говоря, не
постоянен, а коэффициент, определённый через $\operatorname{grad}\mu$,
от состава зависит гораздо слабее (Хартли). Впервые на это указал
Эйнштейн, рассматривая идеальные растворы.
\end{glavnoe}

\begin{primer*}[Опыт Даркена: углерод течёт «в гору»]
Сварим две стали с почти одинаковым содержанием углерода (0,48 и
0,44~\%); первая из них содержит ещё около 4~\% кремния. Кремний повышает
химический потенциал (активность) углерода в железе. После длительного
отжига при 1050\,°C углерод уходит из кремнистой стали в
бескремнистую, хотя там его было \emph{меньше}, и на границе
возникает \emph{скачок} концентрации (рис.~\ref{fig:darken-uphill}):
у шва углерод течёт из области с концентрацией около 0,3~\% в область
с концентрацией около 0,6~\%~— против градиента концентрации, но вниз
по градиенту химического потенциала.
\end{primer*}

\begin{figure}[ht]
\centering
\begin{tikzpicture}
\begin{axis}[width=0.62\textwidth, height=5cm, xmin=-3, xmax=3, ymin=0.2, ymax=0.7,
  xlabel={расстояние от сварного шва}, ylabel={C, мас.\,\%}, xtick={0}, xticklabels={шов},
  ytick={0.3,0.4,0.5,0.6}, yticklabel style={/pgf/number format/use comma}, axis lines=left]
  \addplot[fmGray, dashed, thick] coordinates {(-3,0.48) (0,0.48) (0,0.44) (3,0.44)};
  \addplot[fmRed, thick, smooth] coordinates {(-3,0.48) (-2,0.46) (-1,0.41) (-0.4,0.35) (-0.02,0.31)};
  \addplot[fmRed, thick, smooth] coordinates {(0.02,0.59) (0.4,0.55) (1,0.5) (2,0.455) (3,0.445)};
  \draw[fmRed] (axis cs:0,0.31) -- (axis cs:0,0.59);
  \node[font=\small, anchor=north] at (axis cs:-1.6,0.66) {Fe--C--Si};
  \node[font=\small, anchor=north] at (axis cs:1.6,0.66) {Fe--C};
  \node[font=\scriptsize, fmGray, anchor=north west] at (axis cs:1.2,0.435) {до отжига};
  \node[font=\scriptsize, fmRed, anchor=west] at (axis cs:0.9,0.56) {после отжига};
\end{axis}
\end{tikzpicture}
\caption{«Восходящая» диффузия углерода в опыте Даркена (схема):
углерод уходит из стали с кремнием в сталь без кремния, и у шва
возникает скачок концентрации}
\label{fig:darken-uphill}
\end{figure}

\begin{podrobnee}[сколько независимых коэффициентов]
Коэффициенты $L_{ik}$ не все независимы. При вакансионном механизме
каждому скачку атома отвечает встречный скачок вакансии, и относительно
плоскости, закреплённой в решётке, сумма потоков атомов и вакансий
равна нулю: $\sum_iJ_i=0$. Вместе с соотношениями Онзагера это
оставляет для $n$ компонентов (считая вакансии) $\frac12n(n-1)$
независимых коэффициентов.
\begin{itemize}
\item \emph{Самодиффузия} (атомы + вакансии, $n=2$): $J_1=-J_2=L_{11}(X_1-X_2)$~—
один коэффициент; явление одно, перекрёстных эффектов нет, и
соотношения Онзагера не нужны.
\item \emph{Двойной сплав} (два сорта атомов + вакансии, $n=3$):
\[
J_1=L_{11}(X_1-X_3)+L_{12}(X_2-X_3),\quad
J_2=L_{21}(X_1-X_3)+L_{22}(X_2-X_3),\quad J_3=-(J_1+J_2),
\]
три независимых коэффициента и одно соотношение Онзагера $L_{12}=L_{21}$.
\end{itemize}
\end{podrobnee}
